\documentclass[11pt]{article}
\usepackage[a4paper,margin=27mm]{geometry}
\usepackage{amsmath,amssymb,amsthm,xcolor,booktabs,array}
\newcommand{\PROVED}{\textcolor{green!40!black}{\textbf{PROVED}}}
\newcommand{\OPEN}{\textcolor{red!70!black}{\textbf{OPEN}}}
\newcommand{\AUDIT}{\textcolor{orange!70!black}{\textbf{AUDIT}}}
\newtheorem{theorem}{Theorem}
\newtheorem{corollary}{Corollary}
\title{The Tate moments as the two exterior Gamma states (v122)}
\date{13 August 2026}

\begin{document}
\maketitle

\section{Causal and anticausal boundary states}

For $a>0$ define
\begin{align}
 z_{a,+}(t)&=\sqrt{2a}\int_{-\infty}^t
                  e^{-a(t-s)}f(s)\,ds,\notag\\
 z_{a,-}(t)&=\sqrt{2a}\int_t^\infty
                  e^{-a(s-t)}f(s)\,ds.                       \tag{1}
\end{align}
They satisfy
\begin{equation}
 z_{a,+}'=-az_{a,+}+\sqrt{2a}f,\qquad
 z_{a,-}'=az_{a,-}-\sqrt{2a}f.                              \tag{2}
\end{equation}

\begin{theorem}[Tate moments are exterior Gamma states --- \PROVED]
If $\operatorname{supp}f\subset[L,R]$, then
\begin{align}
 z_{a,+}(t)&=\sqrt{2a}\,e^{-at}
       \int_L^Re^{as}f(s)\,ds &&(t\ge R),\notag\\
 z_{a,-}(t)&=\sqrt{2a}\,e^{at}
       \int_L^Re^{-as}f(s)\,ds &&(t\le L).                   \tag{3}
\end{align}
At $a=\frac12$, the conditions
\begin{equation}
 M_+(f)=\int e^{s/2}f(s)\,ds=0,\qquad
 M_-(f)=\int e^{-s/2}f(s)\,ds=0                             \tag{4}
\end{equation}
are equivalent to vanishing of the right causal and left anticausal
exterior states, respectively.
\end{theorem}

\begin{proof}
For $t\ge R$ and $t\le L$, respectively, pull the $t$-dependent exponential
out of (1).  This gives (3), and (4) is its specialization to $a=1/2$.
\end{proof}

Thus the two-dimensional Tate terminal of v117 is exactly the boundary
state space of the first Gamma cell.  This identification is source-level
and precedes any sign or Schur complement.

\section{A local leakage inequality with a strict margin}

Let $J=[r,r+L]$ and let $z\in H^1(J)$.  For one causal Gamma channel set
\begin{equation}
 f=\frac{z'+az}{\sqrt{2a}},\qquad
 \mathcal E_a(z)=\frac1{a^2}\int_J|z'|^2.                    \tag{5}
\end{equation}

\begin{theorem}[Local Gamma--Tate control --- \PROVED]
For every $a,L>0$,
\begin{equation}
 \boxed{
 \|f\|_{L^2(J)}^2
 \le\left(\frac a2+a^3L^2\right)\mathcal E_a(z)
    +\left(aL+\frac12\right)
       \bigl(|z(r)|^2+|z(r+L)|^2\bigr).}                    \tag{6}
\end{equation}
For the Tate cell $a=\frac12$ this is
\begin{equation}
 \|f\|_{L^2(J)}^2
 \le\left(\frac14+\frac{L^2}{8}\right)\mathcal E_{1/2}(z)
  \frac{L+1}{2}\bigl(|z(r)|^2+|z(r+L)|^2\bigr).             \tag{7}
\end{equation}
In every arithmetic threshold cell,
$L\le\delta_N\le\frac12\log(3/2)<0.203$, so the coefficient of the Gamma
resistance in (7) is smaller than $0.256$.
\end{theorem}

\begin{proof}
Expansion of (5) and integration of the cross term give
\begin{equation}
 \|f\|_2^2
 =\frac1{2a}\|z'\|_2^2+\frac a2\|z\|_2^2
   +\frac12\bigl(|z(r+L)|^2-|z(r)|^2\bigr).                  \tag{8}
\end{equation}
For $t\in J$,
\[
 |z(t)|^2\le2|z(r)|^2+2L\int_J|z'|^2.
\]
The analogous right-endpoint estimate and averaging give
\[
 \|z\|_2^2\le
 L\bigl(|z(r)|^2+|z(r+L)|^2\bigr)+2L^2\|z'\|_2^2.
\]
Substitute this in (8), bound the signed endpoint difference by half the
sum, and use $\|z'\|^2=a^2\mathcal E_a(z)$.  This proves (6).
Equation (7) is $a=1/2$.  Finally $\delta_N$ decreases with $N$ and its
largest value for $N\ge2$ is $\frac12\log(3/2)$.
\end{proof}

\section{Normalization against the moving Tate terminal}

For the old interval $I_T$, the Tate Gram is
\begin{equation}
 G_T=2\begin{pmatrix}\sinh T&T\\T&\sinh T\end{pmatrix}.       \tag{9}
\end{equation}
For a newborn source $e$, write
\[
 b_E=\binom{M_+(e)}{M_-(e)},\qquad
 C=B_T^*G_T^{-1}B_E,\qquad M=I+C^*C,
\]
as in v117.  By (3), the outer $a=1/2$ state vector is
\begin{equation}
 \zeta_E=
 \binom{e^{-R/2}M_+(e)}{e^{L/2}M_-(e)}
 =D_{L,R}b_E,                                               \tag{10}
\end{equation}
with $D_{L,R}=\operatorname{diag}(e^{-R/2},e^{L/2})$.
The Tate terminal energy is
\begin{equation}
 \|CM^{-1/2}e\|^2+\|M^{-1/2}e\|^2=\|e\|^2.                 \tag{11}
\end{equation}

Equations (7), (10), and (11) reduce the clustered-leakage estimate v121
(14) to a finite two-state comparison.  In typed form the candidate is
\begin{equation}
 \boxed{
 \frac{L+1}{2}\,
 B_E^*D_{L,R}^*D_{L,R}B_E
 \le B_E^*G_T^{-1}B_E
   +\left(\frac34-\frac{L^2}{8}\right)I_E.}                  \tag{12}
\end{equation}
Every term in (12) is source-defined.  The scalar remainder is the unused
part of the unit Gamma budget after (7).  Proving (12) uniformly on every
cluster, with the correct endpoint matrices for its orientation, would
complete the unit leakage domination.

\begin{theorem}[Uniform boundary-state domination --- \PROVED]
For every $0<L\le\frac12\log(3/2)$, inequality (12) holds.  In fact its
left side is bounded by $L(1+L)I_E<0.245I_E$, whereas the scalar summand
on the right is larger than $0.744I_E$.
\end{theorem}

\begin{proof}
After translating $J$ to $(0,L)$, the two state functionals in (10) have
representers $e^{(t-L)/2}$ and $e^{-t/2}$.  Their Gram matrix is
\[
 K_L=\begin{pmatrix}
 1-e^{-L}&Le^{-L/2}\\
 Le^{-L/2}&1-e^{-L}
 \end{pmatrix}.
\]
Hence
\[
 \|B_E^*D_{L,R}^*D_{L,R}B_E\|
 =\|K_L\|
 =1-e^{-L}+Le^{-L/2}\le2L.
\]
Multiplication by $(L+1)/2$ bounds the left side of (12) by
$L(1+L)I_E$.  At the largest allowed $L$ this is less than $0.245I_E$.
The scalar part of the right side is
$3/4-L^2/8>0.744$.  The remaining term
$B_E^*G_T^{-1}B_E$ is positive.  This proves (12).
\end{proof}

\begin{center}
\begin{tabular}{>{\raggedright\arraybackslash}p{.55\linewidth}
                >{\centering\arraybackslash}p{.15\linewidth}
                >{\raggedright\arraybackslash}p{.20\linewidth}}
\toprule
Statement & Status & Location \\
\midrule
Tate moments equal the two exterior $a_0$ states & \PROVED & (3)--(4) \\
Strict local Gamma leakage margin & \PROVED & (6)--(7) \\
Explicit state/moment normalization & \PROVED & (9)--(11) \\
Typed boundary-state domination & \PROVED & (12), Theorem 3 \\
Contractivity (8c), condition $G$, row (d) & \OPEN & cascading audit remains \\
\bottomrule
\end{tabular}
\end{center}

\end{document}
